\documentclass[10pt,a4paper]{article}
\usepackage[left=2.5cm,right=2.5cm,top=2.5cm,bottom=2.5cm,]{geometry}
\usepackage{amsmath,amssymb,amsthm}
\newtheorem{theorem}{Theorem}[section]
\newtheorem{lemma}[theorem]{Lemma}
\newtheorem{proposition}[theorem]{Proposition}
\usepackage{setspace}
\setlength{\parindent}{0pt}
\onehalfspacing
\begin{document}
Q: Komari Chika is a thrifty girl who enjoys cooking for herself. 
Vegetable prices at the market fluctuate randomly each day, with a fluctuation range of $\sigma$. 
To prevent financial loss, the magnitude of any price drop is capped at $v$ (where $v < \sigma$). 
Chika does not know the day's price before she goes shopping. 
Suppose that over the past $n$ days, she purchased vegetables $m$ times, resulting in a known price sequence $p = [a_1, \dots, a_m]$. 
She wants to make a purchase at a low price point before the deadline. What's the price can she get when decide to buy vegetables?$\quad$(random walk and optimal stopping? i cant solve now).
\\
\\
Q:Komari Chika has recently become obsessed with the Steam game *Buckshot Roulette*. The rules are as follows: at the start of each round, 
there are $n$ bullets ($n \ge 2$), including at least one blank and one live round. 
The player goes first, followed by the Devil, and they take turns acting. An "action" consists of randomly drawing a bullet; blanks deal no damage, while live rounds cause the loss of one health point. During each action, 
the player can choose to aim the gun at themselves or at the Devil. If the player aims at themselves and draws a blank, the Devil skips their next turn; if they draw a live round, the player loses one health point, 
and the turn proceeds as usual. If the player aims at the Devil and draws a blank, the turn proceeds as usual; if they draw a live round, the Devil loses one health point, and the turn proceeds as usual. 
These rules apply to both parties. The game ends if and only if one party's health drops to zero; if all bullets are exhausted while both parties still have remaining health, the chamber is reloaded.
Assume both parties start with $v$ health points, and the game includes $e$ blanks and $r$ live rounds (where $r + e = n$). Komari Chika wants to determine the strategy that maximizes her win rate.
\\
\\
Q:what is the essence of action? In group theory, consider a act $G\to X$. In fact, it can be seen a homomorphic map $\varphi:G\to\mathrm{Sym}(X)$, 
we can define a equaival relation $\forall a\forall b((a\in G\land b\in G)\rightarrow a\sim b)$, this gives us a split of $G$. we can see the action is a homomorphic map essencally.
But when we consider any object, such give category $\mathcal{C}$, $\forall A\forall B(A\in\mathcal{C}\land B\in\mathcal{C})$, how to understand action?
\\
\\
Q:Show that $\sum\limits_{n=1}^{\infty}\frac{x^n}{1+x+x^2+\cdots+x^{2n-1}}\cos nx$ is uniformly convergent in $[0,1]$.(i am not good at analysis btw.)
\end{document}
